% =====================================================================
%  CHAPITRE 13 — PHRASAL VERBS, IDIOMES, FAUX AMIS
% =====================================================================
\chapter{\en{Phrasal verbs}, expressions idiomatiques et faux amis}

C'est le chapitre qui fait la différence entre un anglais « correct » et un anglais
« vivant ». Les anglophones utilisent des \en{phrasal verbs} en permanence : ils disent
\en{give up}, \en{find out}, \en{turn out}, \en{come up with}, là où un manuel écrit
\en{abandon}, \en{discover}, \en{result}, \en{invent}.

\section{Unité 1 --- Comprendre les \en{phrasal verbs}}

\begin{regle}[Structure]
\textbf{Verbe + particule} (adverbe ou préposition). Le sens est souvent
\textbf{imprévisible} : \en{look} (regarder) + \en{after} (après) = \en{look after}
(s'occuper de). Il faut donc les apprendre \textbf{comme des mots nouveaux}, avec un exemple.
\end{regle}

\begin{tableau}[Les quatre types --- savoir si l'on peut séparer]
\begin{tabularx}{\linewidth}{@{}p{3.2cm}p{4.0cm}X@{}}
\entete{Type} & \entete{Règle} & \entete{Exemples} \tabularnewline
\hline
\textbf{Séparable} & L'objet peut se placer entre le verbe et la particule. Avec un pronom (\en{it, them, him}), la séparation est \textbf{obligatoire}. & \en{turn off the light} = \en{turn the light off} = \en{turn \textbf{it} off} (\en{*turn off it} est faux) \tabularnewline
\textbf{Inséparable} & La particule reste collée au verbe. & \en{look after \textbf{her}} (jamais \en{*look her after}) ; \en{get over it} \tabularnewline
\textbf{Intransitif} & Pas d'objet du tout. & \en{He grew up in Dakar.} \en{The plane took off.} \tabularnewline
\textbf{À trois mots} & Verbe + particule + préposition, toujours inséparable. & \en{look forward to, put up with, get on with, run out of} \tabularnewline
\end{tabularx}
\end{tableau}

\begin{astuce}[La stratégie qui marche : apprendre par particule]
Au lieu d'apprendre des listes alphabétiques, apprenez le \textbf{sens général} de chaque
particule, puis regroupez :\\
\en{up} = achèvement, augmentation (\en{eat up, finish up, speed up, turn up, give up, use up})\\
\en{out} = sortir, apparaître, révéler (\en{go out, find out, work out, run out, take out, look out})\\
\en{off} = séparation, départ, arrêt (\en{take off, turn off, set off, cut off, show off, call off})\\
\en{on} = continuation (\en{go on, carry on, keep on, move on, put on})\\
\en{back} = retour (\en{come back, go back, give back, look back, call back})\\
\en{down} = diminution, arrêt (\en{slow down, calm down, break down, turn down, write down})\\
\en{away} = éloignement, disparition (\en{go away, run away, throw away, put away, take away})\\
\en{over} = au-dessus, réexamen (\en{get over, think over, look over, take over})\\
\en{through} = traversée, achèvement (\en{get through, look through, go through, come through})
\end{astuce}

\subsection{Les 150 \en{phrasal verbs} les plus utiles}
\begin{tableau}[150 \en{phrasal verbs} --- sens et exemple]
\footnotesize
\begin{longtable}{@{}p{3.0cm}p{4.4cm}p{6.4cm}@{}}
\entete{\en{Phrasal verb}} & \entete{Sens français} & \entete{Exemple} \tabularnewline
\endfirsthead
\entete{\en{Phrasal verb}} & \entete{Sens} & \entete{Exemple} \tabularnewline
\endhead
\hline
\en{ask out} & inviter (qqn) à sortir & \en{He asked her out.} \tabularnewline
\en{back up} & sauvegarder ; soutenir & \en{Back up your files. I'll back you up.} \tabularnewline
\en{blow up} & exploser ; gonfler & \en{The car blew up.} \tabularnewline
\en{break down} & tomber en panne ; s'effondrer & \en{My car broke down. She broke down in tears.} \tabularnewline
\en{break into} & s'introduire par effraction & \en{Someone broke into our house.} \tabularnewline
\en{break up} & se séparer ; se terminer & \en{They broke up last month.} \tabularnewline
\en{bring up} & élever ; aborder un sujet & \en{She brought up three kids. Don't bring that up again.} \tabularnewline
\en{build up} & développer progressivement & \en{Build up your vocabulary.} \tabularnewline
\en{call back} & rappeler (téléphone) & \en{I'll call you back.} \tabularnewline
\en{call off} & annuler & \en{They called off the meeting.} \tabularnewline
\en{calm down} & se calmer & \en{Calm down, it's not serious.} \tabularnewline
\en{carry on} & continuer & \en{Carry on working.} \tabularnewline
\en{catch up} & rattraper (son retard) & \en{Let's catch up soon!} \tabularnewline
\en{check in / out} & s'enregistrer / régler sa note & \en{We checked in at the hotel.} \tabularnewline
\en{check out} & examiner ; regarder & \en{Check out this video!} \tabularnewline
\en{cheer up} & remonter le moral & \en{Cheer up, it'll be fine.} \tabularnewline
\en{clean up} & nettoyer, ranger & \en{Clean up your room.} \tabularnewline
\en{come across} & trouver par hasard & \en{I came across an old photo.} \tabularnewline
\en{come from} & venir de & \en{Where do you come from?} \tabularnewline
\en{come in} & entrer & \en{Come in and sit down.} \tabularnewline
\en{come on} & allez ! dépêche-toi & \en{Come on, we're late!} \tabularnewline
\en{come out} & sortir ; paraître & \en{The sun came out. The book comes out in May.} \tabularnewline
\en{come over} & passer chez quelqu'un & \en{Come over for dinner.} \tabularnewline
\en{come up with} & trouver (une idée) & \en{She came up with a great idea.} \tabularnewline
\en{count on} & compter sur & \en{You can count on me.} \tabularnewline
\en{cut down on} & réduire & \en{Cut down on sugar.} \tabularnewline
\en{cut off} & couper ; interrompre & \en{The electricity was cut off. Don't cut me off.} \tabularnewline
\en{deal with} & s'occuper de, traiter & \en{I'll deal with it.} \tabularnewline
\en{depend on} & dépendre de & \en{It depends on the weather.} \tabularnewline
\en{dress up} & s'habiller élégamment & \en{We need to dress up.} \tabularnewline
\en{drop off} & déposer (qqn/qqch) & \en{I'll drop you off at school.} \tabularnewline
\en{eat out} & manger au restaurant & \en{Let's eat out tonight.} \tabularnewline
\en{end up} & finir par & \en{We ended up staying home.} \tabularnewline
\en{fall asleep} & s'endormir & \en{I fell asleep during the film.} \tabularnewline
\en{fall down} & tomber & \en{He fell down the stairs.} \tabularnewline
\en{figure out} & comprendre, résoudre & \en{I can't figure it out.} \tabularnewline
\en{fill in / out} & remplir (un formulaire) & \en{Fill in this form.} \tabularnewline
\en{find out} & découvrir, apprendre & \en{I found out the truth.} \tabularnewline
\en{get along / on with} & bien s'entendre avec & \en{I get along with my colleagues.} \tabularnewline
\en{get away} & s'échapper, partir en vacances & \en{We're getting away for the weekend.} \tabularnewline
\en{get back} & revenir ; récupérer & \en{When did you get back?} \tabularnewline
\en{get in / out of} & entrer / sortir (véhicule) & \en{Get in the car!} \tabularnewline
\en{get off / on} & descendre / monter (bus, train) & \en{Get off at the next stop.} \tabularnewline
\en{get over} & surmonter, se remettre & \en{She got over the illness quickly.} \tabularnewline
\en{get rid of} & se débarrasser de & \en{I got rid of my old phone.} \tabularnewline
\en{get up} & se lever & \en{I get up at six.} \tabularnewline
\en{give away} & donner, offrir ; trahir & \en{Don't give away the ending!} \tabularnewline
\en{give back} & rendre & \en{Give me back my pen.} \tabularnewline
\en{give up} & abandonner ; arrêter & \en{Never give up. He gave up smoking.} \tabularnewline
\en{go ahead} & continuer, vas-y & \en{Go ahead, I'm listening.} \tabularnewline
\en{go back} & retourner & \en{I want to go back home.} \tabularnewline
\en{go on} & continuer ; se passer & \en{Go on! What's going on?} \tabularnewline
\en{go out} & sortir ; s'éteindre & \en{Let's go out. The lights went out.} \tabularnewline
\en{go over} & revoir, vérifier & \en{Let's go over the plan.} \tabularnewline
\en{grow up} & grandir & \en{I grew up in Thiès.} \tabularnewline
\en{hand in} & rendre (un devoir) & \en{Hand in your papers.} \tabularnewline
\en{hang out} & traîner, passer du temps & \en{We hang out at the café.} \tabularnewline
\en{hang up} & raccrocher & \en{Don't hang up!} \tabularnewline
\en{hold on} & attendre, tenir & \en{Hold on a second.} \tabularnewline
\en{hurry up} & se dépêcher & \en{Hurry up or we'll be late.} \tabularnewline
\en{keep on} & continuer & \en{Keep on trying.} \tabularnewline
\en{keep up with} & suivre le rythme & \en{I can't keep up with you.} \tabularnewline
\en{leave out} & omettre, exclure & \en{Don't leave out the details.} \tabularnewline
\en{let down} & décevoir & \en{I won't let you down.} \tabularnewline
\en{listen to} & écouter & \en{Listen to me.} \tabularnewline
\en{live on} & vivre de & \en{They live on rice and fish.} \tabularnewline
\en{look after} & s'occuper de & \en{She looks after her grandmother.} \tabularnewline
\en{look for} & chercher & \en{I'm looking for my keys.} \tabularnewline
\en{look forward to} & attendre avec impatience & \en{I'm looking forward to it.} \tabularnewline
\en{look into} & examiner, enquêter sur & \en{We'll look into the problem.} \tabularnewline
\en{look out} & attention ! & \en{Look out! There's a car!} \tabularnewline
\en{look up} & chercher (un mot) & \en{Look it up in the dictionary.} \tabularnewline
\en{make up} & inventer ; se réconcilier & \en{Don't make up stories. They made up after the fight.} \tabularnewline
\en{move in / out} & emménager / déménager & \en{They moved in last week.} \tabularnewline
\en{pass away} & décéder (poli) & \en{His father passed away in 2020.} \tabularnewline
\en{pay back} & rembourser & \en{I'll pay you back tomorrow.} \tabularnewline
\en{pick up} & ramasser ; aller chercher ; apprendre & \en{Pick it up. I'll pick you up at 8. I picked up some English.} \tabularnewline
\en{point out} & faire remarquer & \en{He pointed out my mistake.} \tabularnewline
\en{put away} & ranger & \en{Put your phone away.} \tabularnewline
\en{put off} & reporter, remettre & \en{Don't put it off until tomorrow.} \tabularnewline
\en{put on} & mettre (vêtement) ; allumer & \en{Put on your coat. Put on the light.} \tabularnewline
\en{put out} & éteindre (feu) & \en{They put out the fire.} \tabularnewline
\en{put up with} & supporter, tolérer & \en{I can't put up with the noise.} \tabularnewline
\en{run into} & rencontrer par hasard & \en{I ran into Ali at the market.} \tabularnewline
\en{run out of} & être à court de & \en{We've run out of sugar.} \tabularnewline
\en{set off} & partir (en voyage) & \en{We set off at dawn.} \tabularnewline
\en{set up} & installer, fonder & \en{They set up a business.} \tabularnewline
\en{show off} & se vanter, frimer & \en{Stop showing off!} \tabularnewline
\en{show up} & arriver, se pointer & \en{He showed up late.} \tabularnewline
\en{shut down} & fermer, éteindre & \en{Shut down your computer.} \tabularnewline
\en{shut up} & se taire (familier, impoli) & \en{Shut up!} --- à éviter avec des inconnus. \tabularnewline
\en{sit down} & s'asseoir & \en{Please sit down.} \tabularnewline
\en{slow down} & ralentir & \en{Slow down, please, I can't follow.} \tabularnewline
\en{speak up} & parler plus fort & \en{Could you speak up?} \tabularnewline
\en{stand up} & se lever & \en{Everybody stood up.} \tabularnewline
\en{switch on / off} & allumer / éteindre & \en{Switch off the TV.} \tabularnewline
\en{take after} & ressembler à (un parent) & \en{She takes after her mother.} \tabularnewline
\en{take away} & emporter & \en{Let's get a takeaway.} \tabularnewline
\en{take care of} & prendre soin de & \en{Take care of yourself.} \tabularnewline
\en{take off} & décoller ; enlever & \en{The plane took off. Take off your shoes.} \tabularnewline
\en{take over} & prendre le contrôle & \en{She took over the company.} \tabularnewline
\en{take up} & commencer (un loisir) ; prendre (de la place) & \en{I took up football. It takes up too much space.} \tabularnewline
\en{talk about} & parler de & \en{Let's talk about your plans.} \tabularnewline
\en{think about / of} & penser à & \en{I'm thinking about moving.} \tabularnewline
\en{think over} & réfléchir à & \en{I'll think it over.} \tabularnewline
\en{throw away} & jeter & \en{Don't throw it away.} \tabularnewline
\en{try on} & essayer (vêtement) & \en{Can I try it on?} \tabularnewline
\en{turn down} & baisser (le son) ; refuser & \en{Turn down the music. He turned down the offer.} \tabularnewline
\en{turn into} & se transformer en & \en{The dream turned into reality.} \tabularnewline
\en{turn off / on} & éteindre / allumer & \en{Turn off your phone.} \tabularnewline
\en{turn out} & s'avérer & \en{It turned out to be true.} \tabularnewline
\en{turn up} & arriver (inattendu) ; monter le son & \en{He turned up at midnight. Turn up the volume.} \tabularnewline
\en{use up} & utiliser jusqu'au bout & \en{We've used up all the money.} \tabularnewline
\en{wake up} & se réveiller & \en{I wake up at five.} \tabularnewline
\en{warm up} & se réchauffer ; s'échauffer & \en{Let's warm up before the match.} \tabularnewline
\en{watch out} & faire attention & \en{Watch out for the dog!} \tabularnewline
\en{wear out} & user ; épuiser & \en{These shoes are worn out. I'm worn out.} \tabularnewline
\en{work out} & fonctionner ; s'entraîner & \en{It worked out fine. I work out twice a week.} \tabularnewline
\en{write down} & noter & \en{Write down the address.} \tabularnewline
\end{longtable}
\end{tableau}

\begin{exo}[1.1 --- Choisissez le \en{phrasal verb}]
\begin{enumerate}
  \item My car (broke down / broke up) on the road.
  \item Please (look after / look for) my bag while I'm away.
  \item I can't (put up with / put on) this noise.
  \item She (takes after / takes off) her father --- they look the same.
  \item We (ran out of / ran into) milk, I need to buy some.
  \item He (made up / made out) an excuse and left.
  \item Can you (speak up / speak out)? I can't hear you.
  \item Don't (give up / give in) --- you're almost there.
  \item I (came across / came up) this photo while cleaning.
  \item They (called off / called back) the match because of the rain.
\end{enumerate}
\end{exo}
\begin{corrige}[1.1]
1. broke down 2. look after 3. put up with 4. takes after 5. ran out of
6. made up 7. speak up 8. give up 9. came across 10. called off
\end{corrige}

\section{Unité 2 --- Les expressions idiomatiques}

Un idiome ne se traduit pas mot à mot : \en{it's raining cats and dogs} ne parle ni de
chats ni de chiens. Apprenez-en \textbf{100} --- cela suffit pour comprendre l'humour,
les films et les conversations normales.

\subsection{Les 40 idiomes les plus fréquents (à connaître absolument)}
\begin{tableau}[40 expressions indispensables]
\footnotesize
\begin{longtable}{@{}p{5.2cm}p{4.8cm}p{3.8cm}@{}}
\entete{Expression} & \entete{Sens français} & \entete{Exemple} \tabularnewline
\endfirsthead
\entete{Expression} & \entete{Sens} & \entete{Exemple} \tabularnewline
\endhead
\hline
\en{a piece of cake} & très facile & \en{The exam was a piece of cake.} \tabularnewline
\en{break a leg!} & bonne chance ! & \en{Break a leg on stage!} \tabularnewline
\en{call it a day} & arrêter pour aujourd'hui & \en{Let's call it a day.} \tabularnewline
\en{cost an arm and a leg} & coûter très cher & \en{That car cost an arm and a leg.} \tabularnewline
\en{don't judge a book by its cover} & ne pas se fier aux apparences & \en{Don't judge a book by its cover.} \tabularnewline
\en{every cloud has a silver lining} & à quelque chose malheur est bon & \en{You lost your job, but every cloud has a silver lining.} \tabularnewline
\en{get cold feet} & se dégonfler & \en{He got cold feet before the wedding.} \tabularnewline
\en{hit the nail on the head} & mettre le doigt dessus & \en{You hit the nail on the head.} \tabularnewline
\en{in the same boat} & dans le même bateau & \en{We're all in the same boat.} \tabularnewline
\en{it's not my cup of tea} & ce n'est pas mon truc & \en{Football is not my cup of tea.} \tabularnewline
\en{keep an eye on} & garder un \oe il sur & \en{Keep an eye on my bag.} \tabularnewline
\en{kill two birds with one stone} & faire d'une pierre deux coups & \en{I killed two birds with one stone.} \tabularnewline
\en{let the cat out of the bag} & vendre la mèche & \en{She let the cat out of the bag.} \tabularnewline
\en{make ends meet} & joindre les deux bouts & \en{It's hard to make ends meet.} \tabularnewline
\en{miss the boat} & rater l'occasion & \en{Don't miss the boat --- apply now.} \tabularnewline
\en{no pain, no gain} & on n'a rien sans rien & \en{No pain, no gain!} \tabularnewline
\en{once in a blue moon} & très rarement & \en{I see him once in a blue moon.} \tabularnewline
\en{on the one hand\ldots on the other hand} & d'un côté\ldots de l'autre & \en{On the one hand, it's cheap\ldots} \tabularnewline
\en{out of the blue} & à l'improviste & \en{He called me out of the blue.} \tabularnewline
\en{pull someone's leg} & faire marcher qqn & \en{Are you pulling my leg?} \tabularnewline
\en{put all your eggs in one basket} & mettre tous ses \oe ufs dans le même panier & \en{Don't put all your eggs in one basket.} \tabularnewline
\en{rain cats and dogs} & pleuvoir des cordes & \en{It's raining cats and dogs.} \tabularnewline
\en{see eye to eye} & être d'accord & \en{We don't see eye to eye on politics.} \tabularnewline
\en{so far, so good} & jusqu'ici, tout va bien & \en{So far, so good.} \tabularnewline
\en{take it easy} & détends-toi, vas-y doucement & \en{Take it easy, you'll be fine.} \tabularnewline
\en{the best of both worlds} & le meilleur des deux mondes & \en{Working from home gives you the best of both worlds.} \tabularnewline
\en{the early bird catches the worm} & le monde appartient à ceux qui se lèvent tôt & \en{Come at 7 --- the early bird catches the worm.} \tabularnewline
\en{there's no such thing as a free lunch} & rien n'est gratuit & \en{There's no such thing as a free lunch.} \tabularnewline
\en{time flies} & le temps passe vite & \en{Time flies when you're busy.} \tabularnewline
\en{to make a long story short} & bref, pour faire court & \en{Long story short, I got the job.} \tabularnewline
\en{under the weather} & pas dans son assiette & \en{I feel a bit under the weather.} \tabularnewline
\en{up in the air} & incertain, en suspens & \en{The plan is still up in the air.} \tabularnewline
\en{when pigs fly} & quand les poules auront des dents & \en{He'll apologise when pigs fly.} \tabularnewline
\en{you can say that again!} & tu l'as dit ! & \en{``It's hot.'' --- ``You can say that again!''} \tabularnewline
\en{your guess is as good as mine} & je n'en sais pas plus que toi & \en{Where is he? Your guess is as good as mine.} \tabularnewline
\en{bite the bullet} & se résigner, s'y mettre & \en{I bit the bullet and called her.} \tabularnewline
\en{break the ice} & briser la glace & \en{A joke helps break the ice.} \tabularnewline
\en{get the hang of it} & prendre le coup de main & \en{You'll get the hang of it.} \tabularnewline
\en{it takes two to tango} & il faut être deux & \en{It takes two to tango.} \tabularnewline
\en{the ball is in your court} & c'est à toi de jouer & \en{I've done my part --- the ball is in your court.} \tabularnewline
\end{longtable}
\end{tableau}

\subsection{80 expressions supplémentaires, par thème}
\begin{tableau}[Expressions de la vie quotidienne]
\footnotesize
\begin{longtable}{@{}p{5.4cm}p{8.4cm}@{}}
\entete{Expression} & \entete{Sens} \tabularnewline
\endfirsthead
\entete{Expression} & \entete{Sens} \tabularnewline
\endhead
\hline
\en{by the way} & au fait, à propos \tabularnewline
\en{that's up to you} & c'est toi qui décides \tabularnewline
\en{it doesn't matter} & ça n'a pas d'importance \tabularnewline
\en{never mind} & laisse tomber, peu importe \tabularnewline
\en{don't worry about it} & ne t'inquiète pas pour ça \tabularnewline
\en{take your time} & prends ton temps \tabularnewline
\en{hurry up, we're late} & dépêche-toi, on est en retard \tabularnewline
\en{I'm on my way} & j'arrive, je suis en route \tabularnewline
\en{I'm running late} & je vais être en retard \tabularnewline
\en{long time no see!} & ça fait longtemps ! \tabularnewline
\en{what's up?} & quoi de neuf ? \tabularnewline
\en{not much / nothing much} & pas grand-chose \tabularnewline
\en{how's it going?} & comment ça va ? \tabularnewline
\en{I can't complain} & je ne peux pas me plaindre \tabularnewline
\en{same old, same old} & toujours la même chose \tabularnewline
\en{I'm stuffed} & je n'en peux plus (de manger) \tabularnewline
\en{I'm starving} & je meurs de faim \tabularnewline
\en{I'm exhausted / knackered} & je suis épuisé \tabularnewline
\en{I'm fed up with it} & j'en ai marre \tabularnewline
\en{it's a pain} & c'est une corvée \tabularnewline
\en{no big deal} & pas grave \tabularnewline
\en{that's a shame / too bad} & c'est dommage \tabularnewline
\en{what a pity!} & quel dommage ! \tabularnewline
\en{congratulations! well done!} & félicitations ! bravo ! \tabularnewline
\en{good luck / all the best} & bonne chance / tous mes v\oe ux \tabularnewline
\en{take care} & prends soin de toi \tabularnewline
\en{see you later / soon} & à plus tard / à bientôt \tabularnewline
\en{have a good one} & bonne journée (AmE) \tabularnewline
\en{my pleasure / you're welcome} & avec plaisir / de rien \tabularnewline
\en{no problem / don't mention it} & pas de souci / il n'y a pas de quoi \tabularnewline
\en{I really appreciate it} & j'apprécie vraiment \tabularnewline
\en{it's my treat} & c'est moi qui invite \tabularnewline
\en{can I get you something?} & je vous sers quelque chose ? \tabularnewline
\en{help yourself} & servez-vous \tabularnewline
\en{make yourself at home} & faites comme chez vous \tabularnewline
\en{I'd better go} & je ferais mieux d'y aller \tabularnewline
\en{I'm afraid I can't} & je crains de ne pas pouvoir \tabularnewline
\en{that's fine by me} & ça me va \tabularnewline
\en{sounds good / sounds great} & ça me paraît bien \tabularnewline
\en{I'm not sure about that} & je n'en suis pas sûr \tabularnewline
\en{I have no idea} & je n'en ai aucune idée \tabularnewline
\en{as far as I know} & pour autant que je sache \tabularnewline
\en{it seems to me that} & il me semble que \tabularnewline
\en{to be honest} & pour être honnête \tabularnewline
\en{frankly speaking} & franchement \tabularnewline
\en{in my opinion / I reckon} & à mon avis / je crois \tabularnewline
\en{I couldn't agree more} & je suis tout à fait d'accord \tabularnewline
\en{I see your point, but} & je comprends votre point de vue, mais \tabularnewline
\en{that's not what I meant} & ce n'est pas ce que je voulais dire \tabularnewline
\en{what do you mean by that?} & qu'entendez-vous par là ? \tabularnewline
\en{you've got a point there} & tu marques un point \tabularnewline
\en{let's agree to disagree} & restons chacun sur nos positions \tabularnewline
\en{it depends on the situation} & cela dépend de la situation \tabularnewline
\en{at the end of the day} & au bout du compte \tabularnewline
\en{the thing is\ldots} & le problème, c'est que\ldots \tabularnewline
\en{to cut a long story short} & pour faire court \tabularnewline
\en{in a nutshell} & en résumé \tabularnewline
\en{that reminds me} & cela me fait penser \tabularnewline
\en{speaking of which} & en parlant de ça \tabularnewline
\en{by any chance} & par hasard \tabularnewline
\en{out of order} & en panne, hors service \tabularnewline
\en{sold out} & épuisé, tout vendu \tabularnewline
\en{on sale} & en solde \tabularnewline
\en{for free / free of charge} & gratuitement \tabularnewline
\en{in advance} & à l'avance \tabularnewline
\en{on time / right on time} & à l'heure / pile à l'heure \tabularnewline
\en{from time to time} & de temps en temps \tabularnewline
\en{all of a sudden} & tout à coup \tabularnewline
\en{as usual} & comme d'habitude \tabularnewline
\en{more or less} & plus ou moins \tabularnewline
\en{pretty much} & à peu près \tabularnewline
\en{kind of / sort of} & un peu, genre \tabularnewline
\en{a lot of / loads of} & beaucoup de \tabularnewline
\en{a bit of a problem} & un petit problème \tabularnewline
\en{the worst part is} & le pire, c'est \tabularnewline
\en{to be on the safe side} & par précaution \tabularnewline
\en{it's worth a try} & ça vaut le coup d'essayer \tabularnewline
\en{give it a go / a shot} & tente le coup \tabularnewline
\en{make the most of it} & profites-en au maximum \tabularnewline
\en{it's now or never} & c'est maintenant ou jamais \tabularnewline
\en{better late than never} & mieux vaut tard que jamais \tabularnewline
\end{longtable}
\end{tableau}

\subsection{L'argot courant (à comprendre, à employer avec modération)}
\begin{tableau}[30 mots familiers très entendus]
\footnotesize
\begin{longtable}{@{}p{3.4cm}p{4.0cm}p{6.4cm}@{}}
\entete{Mot familier} & \entete{Équivalent standard} & \entete{Exemple} \tabularnewline
\endfirsthead
\entete{Familier} & \entete{Standard} & \entete{Exemple} \tabularnewline
\endhead
\hline
\en{guy / dude / mate / bro} & \en{man, friend} & \en{Hey dude, what's up?} \tabularnewline
\en{kids} & \en{children} & \en{She has three kids.} \tabularnewline
\en{stuff / things} & \en{belongings} & \en{Put your stuff here.} \tabularnewline
\en{a lot / tons / loads} & \en{many} & \en{There were loads of people.} \tabularnewline
\en{awesome / amazing / cool} & \en{very good} & \en{That's awesome!} \tabularnewline
\en{awful / terrible / rubbish} & \en{very bad} & \en{The film was rubbish.} \tabularnewline
\en{cheap / pricey} & \en{inexpensive / expensive} & \en{It's a bit pricey.} \tabularnewline
\en{broke} & \en{without money} & \en{I'm broke until payday.} \tabularnewline
\en{bucks / quid} & \en{dollars / pounds} & \en{It costs ten bucks.} \tabularnewline
\en{gig} & \en{job, concert} & \en{I got a new gig.} \tabularnewline
\en{chill / hang out} & \en{relax} & \en{Let's just chill at home.} \tabularnewline
\en{no way!} & \en{impossible!} & \en{No way! You're joking.} \tabularnewline
\en{for real?} & \en{really?} & \en{For real? That's crazy.} \tabularnewline
\en{I'm done} & \en{I have finished} & \en{I'm done with this.} \tabularnewline
\en{my bad} & \en{my mistake, sorry} & \en{My bad, I forgot.} \tabularnewline
\en{no worries} & \en{it's fine} & \en{``Sorry!'' --- ``No worries.''} \tabularnewline
\en{I'm knackered / beat} & \en{I'm exhausted} & \en{I'm beat after work.} \tabularnewline
\en{to be broke / skint} & \en{to have no money} & \en{I'm skint this month.} \tabularnewline
\en{a rip-off} & \en{overpriced} & \en{That price is a rip-off.} \tabularnewline
\en{to be into something} & \en{to like} & \en{I'm really into football.} \tabularnewline
\en{to ghost someone} & \en{to stop replying} & \en{He ghosted me after two dates.} \tabularnewline
\en{to flake} & \en{to cancel at the last minute} & \en{Don't flake on me.} \tabularnewline
\en{to screw up} & \en{to make a mistake} & \en{I screwed up the interview.} \tabularnewline
\en{to be fed up} & \en{to be annoyed} & \en{I'm fed up with waiting.} \tabularnewline
\en{to hang in there} & \en{to persevere} & \en{Hang in there, it'll get better.} \tabularnewline
\en{to play it by ear} & \en{to improvise} & \en{Let's play it by ear.} \tabularnewline
\en{to be on the same page} & \en{to agree} & \en{Are we on the same page?} \tabularnewline
\en{to hit the sack} & \en{to go to bed} & \en{I'm hitting the sack.} \tabularnewline
\en{what's the catch?} & \en{what's the hidden problem?} & \en{Free? What's the catch?} \tabularnewline
\en{it's a no-brainer} & \en{it's an obvious choice} & \en{Take the job --- it's a no-brainer.} \tabularnewline
\end{longtable}
\end{tableau}

\section{Unité 3 --- Les faux amis : 150 mots qui vous trahissent}

Un faux ami est un mot qui ressemble au français mais ne veut pas dire la même chose.
Ce sont les fautes les plus dangereuses : personne ne vous corrige, on vous comprend mal
--- et vous ne le savez pas.

\begin{tableau}[150 faux amis français \ra anglais]
\footnotesize
\begin{longtable}{@{}p{3.0cm}p{3.6cm}p{3.6cm}p{3.6cm}@{}}
\entete{Mot français} & \entete{Ce que dit l'anglais} & \entete{Le mot anglais qui ressemble} & \entete{Ce qu'il veut dire vraiment} \tabularnewline
\endfirsthead
\entete{Français} & \entete{Anglais correct} & \entete{Faux ami} & \entete{Sens réel} \tabularnewline
\endhead
\hline
actuellement & \en{currently, at present} & \en{actually} & en fait, en réalité \tabularnewline
éventuellement & \en{possibly} & \en{eventually} & finalement, à la fin \tabularnewline
attendre & \en{to wait (for)} & \en{to attend} & assister à, être présent \tabularnewline
assister à & \en{to attend} & \en{to assist} & aider \tabularnewline
blesser & \en{to injure, hurt} & \en{to bless} & bénir \tabularnewline
librairie & \en{bookshop} & \en{library} & bibliothèque \tabularnewline
bibliothèque & \en{library} & \en{librairie} (rare) & --- \tabularnewline
blessé & \en{injured} & \en{blessed} & béni \tabularnewline
caméra & \en{video camera} & \en{camera} & appareil photo \tabularnewline
appareil photo & \en{camera} & --- & --- \tabularnewline
monnaie (pièces) & \en{change, coins} & \en{money} & argent en général \tabularnewline
monnaie (devise) & \en{currency} & \en{money} & argent \tabularnewline
déjeuner & \en{lunch} & \en{dinner} & dîner \tabularnewline
dîner & \en{dinner} & \en{diner} & personne qui dîne ; petit restaurant \tabularnewline
pain & \en{bread} & \en{pain} & douleur \tabularnewline
travail & \en{work, job} & \en{travail} (rare) & --- \tabularnewline
travailler & \en{to work} & \en{to travel} & voyager \tabularnewline
voyage & \en{trip, journey} & \en{voyage} & traversée maritime \tabularnewline
magasin & \en{shop, store} & \en{magazine} & revue \tabularnewline
revue & \en{magazine} & --- & --- \tabularnewline
blessure & \en{wound, injury} & --- & --- \tabularnewline
envie & \en{desire, wish} & \en{envy} & jalousie \tabularnewline
jalousie & \en{jealousy} & --- & --- \tabularnewline
sensible & \en{sensitive} & \en{sensible} & raisonnable, sensé \tabularnewline
raisonnable & \en{reasonable} & \en{sensible} & sensé \tabularnewline
coin & \en{corner} & \en{coin} & pièce de monnaie \tabularnewline
pièce & \en{room, coin} & --- & --- \tabularnewline
chance & \en{luck} & \en{chance} & occasion, hasard \tabularnewline
hasard & \en{chance, randomness} & \en{hazard} & danger \tabularnewline
danger & \en{danger} & \en{hazard} & risque \tabularnewline
médecine & \en{medicine} & \en{medication} & traitement médicamenteux \tabularnewline
médicament & \en{medicine, drug, pill} & \en{medicine} & aussi : la médecine \tabularnewline
drogue & \en{drug} & \en{drug} & médicament ET drogue \tabularnewline
pharmacie & \en{pharmacy, chemist's} & \en{pharmacy} & idem \tabularnewline
préservatif & \en{condom} & \en{preservative} & conservateur (alimentaire) \tabularnewline
conservateur & \en{preservative, conservative} & \en{preservative} & additif alimentaire \tabularnewline
expérience & \en{experience} & \en{experiment} & expérience scientifique \tabularnewline
expérience (labo) & \en{experiment} & --- & --- \tabularnewline
essayer & \en{to try} & \en{to essay} (rare) & --- \tabularnewline
essai & \en{attempt, essay} & \en{essay} & dissertation \tabularnewline
sujet & \en{subject, topic} & \en{subject} & aussi : sujet grammatical \tabularnewline
objet & \en{object, item} & \en{objective} & objectif \tabularnewline
objectif & \en{goal, objective} & \en{objective} & objectif (impartial) \tabularnewline
réellement & \en{really} & \en{actually} & en fait \tabularnewline
actuel & \en{current, present} & \en{actual} & réel \tabularnewline
issue & \en{exit} & \en{issue} & problème, question \tabularnewline
sortie & \en{exit, way out} & --- & --- \tabularnewline
blessé & \en{injured} & --- & --- \tabularnewline
rester & \en{to stay} & \en{to rest} & se reposer \tabularnewline
se reposer & \en{to rest} & --- & --- \tabularnewline
demander & \en{to ask} & \en{to demand} & exiger \tabularnewline
exiger & \en{to demand, require} & --- & --- \tabularnewline
répondre & \en{to answer, reply} & \en{to respond} & réagir \tabularnewline
proposer & \en{to suggest, offer} & \en{to propose} & demander en mariage ; proposer un plan \tabularnewline
soutenir & \en{to support} & \en{to sustain} & maintenir \tabularnewline
supporter (tolérer) & \en{to stand, put up with} & \en{to support} & soutenir \tabularnewline
introduire & \en{to introduce (people)} & \en{to introduce} & présenter quelqu'un \tabularnewline
présenter & \en{to introduce, present} & --- & --- \tabularnewline
connaître & \en{to know} & \en{to recognise} & reconnaître \tabularnewline
reconnaître & \en{to recognise} & --- & --- \tabularnewline
manquer & \en{to miss} & \en{to lack} & manquer de, être dépourvu \tabularnewline
apprendre & \en{to learn} & \en{to teach} & enseigner \tabularnewline
enseigner & \en{to teach} & --- & --- \tabularnewline
emprunter & \en{to borrow} & \en{to lend} & prêter \tabularnewline
prêter & \en{to lend} & \en{to borrow} & emprunter \tabularnewline
apporter & \en{to bring} & \en{to take} & emporter \tabularnewline
emporter & \en{to take (away)} & \en{to bring} & apporter \tabularnewline
visiter (un lieu) & \en{to visit} & \en{to visit} & idem (mais on ne « visite » pas une personne : \en{to visit somebody} existe, \en{to see somebody} est plus naturel) \tabularnewline
rendre visite & \en{to visit, to see} & --- & --- \tabularnewline
occasion & \en{opportunity, chance} & \en{occasion} & événement, circonstance \tabularnewline
parfois & \en{sometimes} & \en{often} & souvent \tabularnewline
souvent & \en{often} & --- & --- \tabularnewline
peut-être & \en{maybe, perhaps} & \en{possibly} & éventuellement \tabularnewline
possible & \en{possible} & \en{feasible} & réalisable \tabularnewline
important & \en{important} & \en{impressive} & impressionnant \tabularnewline
impressionnant & \en{impressive} & --- & --- \tabularnewline
intéressant & \en{interesting} & \en{interested} & intéressé \tabularnewline
intéressé & \en{interested} & --- & --- \tabularnewline
ennuyeux & \en{boring} & \en{bored} & ennuyé \tabularnewline
ennuyé & \en{annoyed, bothered} & \en{bored} & qui s'ennuie \tabularnewline
fatigué & \en{tired} & \en{fatigued} & épuisé (médical) \tabularnewline
malade & \en{ill, sick} & \en{malade} (inexistant) & --- \tabularnewline
médecin & \en{doctor} & \en{medical} & médical \tabularnewline
hôpital & \en{hospital} & \en{hotel} & hôtel \tabularnewline
hôtel & \en{hotel} & \en{hospital} & hôpital \tabularnewline
hôpital/hôtel & --- & \en{hostel} & auberge \tabularnewline
monnaie/pièce & \en{coin} & --- & --- \tabularnewline
billet & \en{ticket, banknote} & \en{bill} & facture, addition (AmE) \tabularnewline
facture & \en{invoice, bill} & --- & --- \tabularnewline
recette & \en{recipe} & \en{receipt} & reçu, ticket de caisse \tabularnewline
reçu & \en{receipt} & \en{recipe} & recette de cuisine \tabularnewline
prévenir & \en{to warn} & \en{to prevent} & empêcher \tabularnewline
empêcher & \en{to prevent, stop} & --- & --- \tabularnewline
réussir & \en{to succeed} & \en{to success} (n'existe pas) & --- \tabularnewline
réussir un examen & \en{to pass} & \en{to succeed} & réussir (dans la vie) \tabularnewline
échouer & \en{to fail} & \en{to fall} & tomber \tabularnewline
tomber & \en{to fall} & --- & --- \tabularnewline
monter & \en{to go up, climb} & \en{to mount} & assembler, monter (technique) \tabularnewline
descendre & \en{to go down} & \en{to descend} & descendre (littéraire) \tabularnewline
partir & \en{to leave} & \en{to part} & se séparer \tabularnewline
quitter & \en{to leave} & \en{to quit} & arrêter, démissionner \tabularnewline
arrêter & \en{to stop} & \en{to arrest} & arrêter (police) \tabularnewline
se marier & \en{to get married} & \en{to marry} & épouser \tabularnewline
divorcer & \en{to get divorced} & \en{to divorce} & idem \tabularnewline
enceinte & \en{pregnant} & \en{enceinte} (n'existe pas) & --- \tabularnewline
grossesse & \en{pregnancy} & --- & --- \tabularnewline
accoucher & \en{to give birth} & \en{to accouche} (n'existe pas) & --- \tabularnewline
couche & \en{diaper / nappy} & --- & --- \tabularnewline
couche (strate) & \en{layer} & --- & --- \tabularnewline
élever (un enfant) & \en{to raise, bring up} & \en{to elevate} & soulever, hausser \tabularnewline
monter (prix) & \en{to rise, increase} & \en{to raise} & augmenter (qqch) \tabularnewline
augmenter & \en{to increase} & \en{to augment} (rare) & --- \tabularnewline
diminuer & \en{to decrease} & \en{to diminish} & diminuer (soutenu) \tabularnewline
changer & \en{to change} & \en{to exchange} & échanger \tabularnewline
échanger & \en{to exchange} & --- & --- \tabularnewline
échanger (monnaie) & \en{to change, exchange} & --- & --- \tabularnewline
dépenser & \en{to spend} & \en{to dispense} & distribuer \tabularnewline
dépens & --- & \en{expense} & frais, dépense \tabularnewline
coûter & \en{to cost} & \en{to count} & compter \tabularnewline
compter & \en{to count} & \en{to compute} & calculer \tabularnewline
payer & \en{to pay} & \en{to pay for} & payer pour (au sens figuré) \tabularnewline
vendre & \en{to sell} & \en{to sale} (n'existe pas) & --- \tabularnewline
vente & \en{sale} & --- & --- \tabularnewline
acheter & \en{to buy} & \en{to book} & réserver \tabularnewline
réserver & \en{to book, reserve} & --- & --- \tabularnewline
livrer & \en{to deliver} & \en{to delivery} (n'existe pas) & --- \tabularnewline
livraison & \en{delivery} & --- & --- \tabularnewline
envoyer & \en{to send} & \en{to envy} & envier \tabularnewline
envier & \en{to envy} & --- & --- \tabularnewline
recevoir & \en{to receive} & \en{to perceive} & percevoir \tabularnewline
apercevoir & \en{to notice, glimpse} & --- & --- \tabularnewline
écouter & \en{to listen to} & \en{to hear} & entendre \tabularnewline
entendre & \en{to hear} & --- & --- \tabularnewline
regarder & \en{to watch, look at} & \en{to see} & voir \tabularnewline
voir & \en{to see} & --- & --- \tabularnewline
sentir (odeur) & \en{to smell} & \en{to feel} & ressentir, toucher \tabularnewline
sentir (ressentir) & \en{to feel} & \en{to sense} & pressentir \tabularnewline
toucher & \en{to touch} & \en{to feel} & palper, ressentir \tabularnewline
goûter (n.) & \en{snack} & \en{taste} & goût \tabularnewline
goûter (v.) & \en{to taste} & --- & --- \tabularnewline
savourer & \en{to savour, enjoy} & --- & --- \tabularnewline
assaisonnement & \en{seasoning} & \en{season} & saison \tabularnewline
saison & \en{season} & --- & --- \tabularnewline
temps (météo) & \en{weather} & \en{time} & temps (durée) \tabularnewline
heure & \en{hour, time} & --- & --- \tabularnewline
année & \en{year} & \en{yearly} & annuel \tabularnewline
journée & \en{day} & \en{journey} & voyage \tabularnewline
voyage & \en{trip, journey} & --- & --- \tabularnewline
promenade & \en{walk} & \en{promenade} & front de mer \tabularnewline
aéroport & \en{airport} & \en{port} & port maritime \tabularnewline
port & \en{port, harbour} & --- & --- \tabularnewline
aéroport & --- & \en{airfield} & aérodrome \tabularnewline
voiture & \en{car} & \en{vehicle} & véhicule \tabularnewline
camion & \en{truck / lorry} & \en{camera} & appareil photo \tabularnewline
car (bus) & \en{coach} & \en{car} & voiture \tabularnewline
train & \en{train} & \en{tram} & tramway \tabularnewline
tramway & \en{tram / streetcar} & --- & --- \tabularnewline
métro & \en{underground / subway / metro} & \en{metro} & idem (AmE : \en{subway}) \tabularnewline
billet & \en{ticket} & \en{billet} (rare) & --- \tabularnewline
rue & \en{street} & \en{road} & route \tabularnewline
route & \en{road} & \en{route} & itinéraire \tabularnewline
itinéraire & \en{route, itinerary} & --- & --- \tabularnewline
adresse & \en{address} & \en{advice} & conseil \tabularnewline
conseil & \en{advice} & \en{advert} & publicité \tabularnewline
publicité & \en{advertising, advert} & \en{publicity} & notoriété \tabularnewline
réclame & --- & \en{publicity} & notoriété, battage \tabularnewline
information & \en{information} & \en{inform} & informer \tabularnewline
nouvelles & \en{news} & --- & --- \tabularnewline
journal & \en{newspaper} & \en{journal} & revue spécialisée ; journal intime \tabularnewline
journal intime & \en{diary} & \en{diary} & agenda (BrE) \tabularnewline
agenda & \en{diary / planner} & \en{agenda} & ordre du jour \tabularnewline
ordre du jour & \en{agenda} & --- & --- \tabularnewline
réunion & \en{meeting} & \en{union} & syndicat, union \tabularnewline
syndicat & \en{union, trade union} & --- & --- \tabularnewline
bureau & \en{desk / office} & \en{bureau} & agence, service \tabularnewline
agence & \en{agency} & --- & --- \tabularnewline
entreprise & \en{company, business} & \en{enterprise} & esprit d'entreprise \tabularnewline
usine & \en{factory, plant} & \en{enterprise} & --- \tabularnewline
patron & \en{boss} & \en{patron} & mécène, client régulier \tabularnewline
client & \en{customer, client} & --- & --- \tabularnewline
personne (individu) & \en{person} & \en{people} & les gens \tabularnewline
gens & \en{people} & \en{persons} & personnes (formel) \tabularnewline
homme & \en{man} & \en{human} & humain \tabularnewline
humain & \en{human} & --- & --- \tabularnewline
monsieur & \en{Mr / sir} & \en{monsieur} (n'existe pas) & --- \tabularnewline
madame & \en{Mrs / Ms / madam} & --- & --- \tabularnewline
demoiselle & \en{Miss} & --- & --- \tabularnewline
nom & \en{name} & \en{noun} & nom (grammaire) \tabularnewline
verbe & \en{verb} & \en{adverb} & adverbe \tabularnewline
phrase & \en{sentence} & \en{phrase} & expression, groupe de mots \tabularnewline
peine (sanction) & \en{sentence} & \en{penalty} & pénalité \tabularnewline
erreur & \en{mistake, error} & \en{error} & erreur (technique) \tabularnewline
faute & \en{mistake, fault} & \en{fault} & responsabilité, défaut \tabularnewline
défaut & \en{flaw, fault} & \en{default} & défaut (informatique) \tabularnewline
succès & \en{success} & \en{successful} & réussi \tabularnewline
chance & \en{luck, opportunity} & \en{lucky} & chanceux \tabularnewline
malchance & \en{bad luck} & \en{unlucky} & malchanceux \tabularnewline
\end{longtable}
\end{tableau}

\begin{piege}[Les cinq faux amis les plus coûteux]
\begin{enumerate}
  \item \en{actually} = \textbf{en fait} (pas « actuellement » \ra \en{currently}).
  \item \en{eventually} = \textbf{finalement} (pas « éventuellement » \ra \en{possibly}).
  \item \en{library} = \textbf{bibliothèque} (pas « librairie » \ra \en{bookshop}).
  \item \en{to attend} = \textbf{assister à} (pas « attendre » \ra \en{to wait}).
  \item \en{sensible} = \textbf{raisonnable} (pas « sensible » \ra \en{sensitive}).
\end{enumerate}
\end{piege}

\section{Unité 4 --- Les calques du français à éviter}
\begin{tableau}[Traductions littérales fautives et formes correctes]
\footnotesize
\begin{longtable}{@{}p{5.4cm}p{4.0cm}p{4.4cm}@{}}
\entete{\ko Calque du français} & \entete{\ok Anglais correct} & \entete{Français} \tabularnewline
\endfirsthead
\entete{\ko Calque} & \entete{\ok Correct} & \entete{Français} \tabularnewline
\endhead
\hline
\en{*I have 30 years.} & \en{I'm 30 (years old).} & J'ai 30 ans. \tabularnewline
\en{*I am agree.} & \en{I agree.} & Je suis d'accord. \tabularnewline
\en{*I have cold / hunger.} & \en{I'm cold / hungry.} & J'ai froid / faim. \tabularnewline
\en{*It depends of you.} & \en{It depends on you.} & Cela dépend de toi. \tabularnewline
\en{*I listen music.} & \en{I listen to music.} & J'écoute de la musique. \tabularnewline
\en{*I search my keys.} & \en{I'm looking for my keys.} & Je cherche mes clés. \tabularnewline
\en{*I explain you the problem.} & \en{I'll explain the problem to you.} & Je t'explique le problème. \tabularnewline
\en{*I said you something.} & \en{I told you something.} & Je t'ai dit quelque chose. \tabularnewline
\en{*He is married with a doctor.} & \en{He is married to a doctor.} & Il est marié à une médecin. \tabularnewline
\en{*I will see you to tomorrow.} & \en{I'll see you tomorrow.} & Je te verrai demain. \tabularnewline
\en{*In the morning I go in the office.} & \en{In the morning I go to the office.} & Le matin je vais au bureau. \tabularnewline
\en{*I am in Dakar since 2015.} & \en{I have been in Dakar since 2015.} & Je suis à Dakar depuis 2015. \tabularnewline
\en{*It is difficult for me of understand.} & \en{It is difficult for me to understand.} & C'est difficile pour moi de comprendre. \tabularnewline
\en{*I have fear.} & \en{I'm afraid.} & J'ai peur. \tabularnewline
\en{*She makes a mistake.} \emph{(correct, mais)} & \en{She made a mistake.} (passé simple, pas présent) & Elle a fait une erreur. \tabularnewline
\en{*I demand you a question.} & \en{I'll ask you a question.} & Je te pose une question. \tabularnewline
\en{*This movie passes on TV.} & \en{This movie is on TV.} & Ce film passe à la télé. \tabularnewline
\en{*I have the possibility of coming.} & \en{I can come.} / \en{I have the chance to come.} & J'ai la possibilité de venir. \tabularnewline
\en{*Make attention!} & \en{Be careful!} / \en{Watch out!} & Fais attention ! \tabularnewline
\en{*I am a Senegalese.} & \en{I'm Senegalese.} / \en{I'm from Senegal.} & Je suis sénégalais. \tabularnewline
\en{*It is not possible of doing that.} & \en{It isn't possible to do that.} & Ce n'est pas possible de faire ça. \tabularnewline
\en{*I wait your answer.} & \en{I'm waiting for your answer.} / \en{I look forward to your reply.} & J'attends votre réponse. \tabularnewline
\en{*The weather is very hot.} \emph{(correct)} & \en{It's very hot.} (plus naturel) & Il fait très chaud. \tabularnewline
\en{*I want that you come.} & \en{I want you to come.} & Je veux que tu viennes. \tabularnewline
\en{*I think that it is true.} \emph{(correct)} & \en{I think it's true.} (\en{that} omis) & Je pense que c'est vrai. \tabularnewline
\end{longtable}
\end{tableau}

\begin{mission}[Chapitre 13 --- 21 jours d'expressions]
\textbf{Semaine 1 :} 30 \en{phrasal verbs} (les plus fréquents de la liste) --- 5 par jour,
chacun dans une phrase personnelle, à voix haute.\\
\textbf{Semaine 2 :} 25 idiomes --- 5 par jour, chacun dans une mini-situation jouée à voix haute.\\
\textbf{Semaine 3 :} 25 faux amis --- 5 par jour, avec la phrase correcte.\\
À la fin, vous emploierez \en{figure out, get rid of, look forward to, it depends on,
make ends meet, by the way} --- et votre anglais cessera de sonner traduit.
\end{mission}
